\documentclass[12pt]{article}
\usepackage[margin=1in]{geometry}
\usepackage{amsmath}
\usepackage{fancyhdr}
\fancypagestyle{plain}{%
  \fancyhf{}
  \fancyhead[L]{POLISCI 450A}
  \fancyhead[R]{Fall 2026}
  %\fancyfoot[C]{\thepage}
}
\newcommand{\var}{\operatorname{Var}}
\newcommand{\cov}{\operatorname{Cov}}
\newcommand{\E}{\operatorname{\mathbb{E}}}
\newcommand{\solution}{\bigskip\textbf{\textit{Solution:}}\xspace}
\DeclareMathOperator*\plim{plim}

\usepackage{titlesec}
\titleformat{\section}{\large\bfseries}{\thesection}{1em}{}
\usepackage{titling}
\setlength{\droptitle}{-3em}   % move title up


\title{\Large Lecture 5-1: Regression with Two Regressors}
\author{}
\date{}
\begin{document}
\maketitle
\vspace{-6em}

\noindent \textbf{Goal:} Introduce mechanics of multiple regression with two regressors.\medskip

\begin{itemize}\setlength{\itemsep}{1em}
\item Motivation for multiple regression
\begin{itemize}
\item Richer description and prediction
\item Control for confounders
\item Model nonlinearity and interactions
\item Parameters, estimators, estimates
\end{itemize}
\item Two cases
\begin{itemize}
\item Binary + continuous
\item Continuous + continuous
\end{itemize}
\item Mechanics of multiple regression
\begin{itemize}
\item Derivation and Properties of OLS
\item Analysis of Variance
\item Partialling out interpretation
\end{itemize}
\end{itemize}


% End outline
\clearpage


\section*{Worksheet}

\begin{enumerate}\itemsep3em

\item (\emph{Interpreting a Binary Covariate in a Two–Regressor Model})  
Suppose
\[
\widehat{y}=\hat\beta_0+\hat\beta_1 x_1+\hat\beta_2 x_2,\qquad
(\hat\beta_0,\hat\beta_1,\hat\beta_2)=(78.3,\,-10.6,\,-21.5),
\]
where $y$ is executive corruption (0--100), $x_1$ is high court independence (0--4), and $x_2\in\{0,1\}$ indicates democracy.
\begin{enumerate}
    \item For a fixed level of court independence $x_1$, compute the predicted difference $\widehat{y}(x_2=1)-\widehat{y}(x_2=0)$ and interpret.  
    \item For a fixed regime type, compute the predicted change when $x_1$ increases by one unit.  
    \item Verify that the slope on court independence is the same in both groups and explain why.
\end{enumerate}

\item (\emph{Regression Plane and Counterfactual Prediction})  
With two continuous covariates, suppose
\[
\widehat{y}=\hat\beta_0+\hat\beta_1 x_1+\hat\beta_2 x_2,\qquad
(\hat\beta_0,\hat\beta_1,\hat\beta_2)=(166.4,\,-13.95,\,-22.0),
\]
where $y$ is executive corruption (0--100), $x_1$ is high court independence (0--4), and $x_2$ is $\log_{10}$ GDP per capita.
(a) Show that for any two units $A$ and $B$,
\[
\widehat{y}_A-\widehat{y}_B=\hat\beta_1(x_{1A}-x_{1B})+\hat\beta_2(x_{2A}-x_{2B}).
\]
(b) Using $(x_1,x_2)=(2.96,\,4.48)$ for Chile and $(0.21,\,4.38)$ for China, compute the predicted levels and their difference.  
(c) Explain the ``partial effect'' meaning of $\hat\beta_1$ and $\hat\beta_2$.

\item (\emph{Properties of the fitted values and residuals}) Show that
\begin{enumerate}
  \item $\sum_{i=1}^n \hat{u}_i=0 \Rightarrow \bar{y}=\hat{\bar{y}}$.
\item The point ($\bar{y}, \bar{x},\bar{z}$) is always on the regression plane.
\item $Cov(x,\hat{u})=Cov(z,\hat{u})=0 \Rightarrow Cov(\hat{y},\hat{u})=0$.
\end{enumerate}


\item (\emph{ANOVA Decomposition, $R^2$, and Adjusted $R^2$})  
Let $\hat y_i$ be fitted values from OLS with intercept and two regressors, and $\hat u_i=y_i-\hat y_i$. Prove that
\[
\sum_{i=1}^n(y_i-\bar y)^2=\sum_{i=1}^n(\hat y_i-\bar y)^2+\sum_{i=1}^n\hat u_i^2\quad(\text{SST}=\text{SSE}+\text{SSR}).
\]
Then explain why $R^2=\text{SSE}/\text{SST}$ is weakly increasing when a regressor is added.


\end{enumerate}


\end{document}
